% Capítulo 6 (Resultados) — Slide 11 — CTI ajustado: leitura
\begin{frame}{CTI ajustado: leitura (análise complementar)}
  \[
    J(\pi,g) = \mathrm{CTI}_{\text{norm}}(\pi,g) + \lambda_{BE}\cdot \mathrm{BE}_{\text{norm}}(\pi,g)
  \]
  \begin{itemize}
    \small
    \item $\lambda_{BE}=0$: recupera o CTI puro (EOQ domina)
    \item $\lambda_{BE}\geq 0{,}25$: política única migra para
      \highlight{SA} (BE\,$=$\,5,5 vs.\ 55,4 do EOQ)
    \item Em $\lambda_{BE}=1{,}0$: até +56,6\% de redução ajustada, com NS
      caindo de 0,942 para 0,862
  \end{itemize}
  \vfill
  \begin{alertblock}{Status: preliminar e complementar}
    \footnotesize Não substitui o critério principal (CTI puro sob
    NS $\geq 0{,}70$). A partir de $\lambda_{BE}\geq 0{,}5$, a seleção por
    perfil converge para a política única ajustada.
  \end{alertblock}
  \note{
    É fundamental destacar que essa é uma análise preliminar e complementar
    -- o critério principal da dissertação continua sendo o CTI puro sob a
    restrição de NS mínimo. Notem também que, quando a penalidade por
    instabilidade domina o critério, a seleção por perfil perde sua
    vantagem diferenciadora, porque a política mais estável passa a ser
    preferida uniformemente em todos os perfis, convergindo para o mesmo
    resultado da política única ajustada.
  }
\end{frame}
